\chapter{证据索引}
\renewcommand{\thesubsection}{\thesection.\arabic{subsection}}

2026-10-31 核验：远端 \texttt{c385a9d99754} 的板测更新已合并到工作区；本稿最终演示对象改为 PUF3。主要数据来自 \pufresult，工件位于 \pufcandidate。旧 PYNQ release、PL-only 与当前 PUF3 按原身份引用，不能把旧版已通过项目自动转移到当前版本。

本目录的 \texttt{review\_evidence.py} 从原始 JSON、manifest 与报告生成下列摘录，并验证哈希、周期重算与正文主表。输出 \texttt{validation/evidence\_check.json} 记录证据文件哈希。墙钟在 PUF3 JSON 中只有摘要，不能独立复算其逐次中位数。ISE 配图对应同日另列的 XSim 单元复验；报告构建仅读取证据，没有重新综合实现或上板。

{\small
\setlength{\tabcolsep}{4pt}
\renewcommand{\thetable}{C-1}
\begin{longtable}{@{}>{\raggedright\arraybackslash}p{0.31\linewidth}>{\raggedright\arraybackslash}p{0.62\linewidth}@{}}
\caption{结论与主要证据对照}\label{tab:C-1}\\
\toprule
\textbf{结论} & \textbf{主要证据} \\
\midrule
\endfirsthead
\multicolumn{2}{@{}l}{\footnotesize 续表 C-1}\\
\toprule
\textbf{结论} & \textbf{主要证据} \\
\midrule
\endhead
当前工件身份 & \pufcandidate 的 manifest.json、sha256sums.txt \\
当前实现/资源/时序 & \pufreports 的 timing\_summary.rpt、utilization.rpt、drc.rpt、ro\_corner\_sta.txt \\
当前 KAT/benchmark/PUF & \pufresult 及同目录 mailbox\_log.txt \\
当前执行 Notebook & \texttt{src/pynq\_puf3/pynq\_puf3\_board\_test.ipynb} \\
当前仿真与固件 & candidate/reports/cosim\_summary.json、firmware\_build.json \\
旧 release 全容量统计 & \texttt{data/pynq/board\_20261028/pynq\_board\_full.json}，只用于历史表 2-16 \\
旧 release 错误恢复 & \texttt{data/pynq/error\_gate\_20261031\_112812/}；执行本体在 src/pynq/ \\
host 标准与失败契约 & \texttt{data/ascon/results/} 的 host\_validation、reference\_gate、failure\_contract JSON \\
历史 PL-only 原始日志 & \texttt{data/ascon/results/board/} 下 uart\_20261026\_sw\_and\_ise.txt、uart\_20261027\_mmio.txt、uart\_20261027\_keys3.txt \\
历史策略记录 & \texttt{data/ascon/results/strategy\_sweep\_20261017/} \\
真实工程截图 & 本报告 figs/ 下 Block Design、Implemented Design、Utilization、ISE waveform 四张 PNG；用户提供，原字节保存 \\
ISE 截图配套自检 & \texttt{data/ascon/results/\allowbreak{}ise\_screenshot\_20261031/}，1502 checks / 0 errors \\
实物照片 & \texttt{docs/tech\_report/figs/board\_photo.jpg}；不证明当前运行镜像或芯片速度等级 \\
\bottomrule
\end{longtable}}

资源利用率界面和 ISE 真实仿真波形已补入图 3-2、3-5，原工程截图的缺图位置已补齐。截图支持可视核对，不替代原始报告、板测身份或 RO 全域签核。历史本地备份不保证随 clone 可用；最终证据包须按此稿与新数据重新装配。

\Needspace{85pt}
\section{当前 PUF3 工件身份}
来源：\pufcandidate 的 manifest.json。以下仅摘取身份与工件字段，构建时的 NOT\_RUN 保持原样，后续板测另列。
\begin{codeblock}
{
  "source_commit": "27f24bb072f53c3eb109a805ed1ecba3869c0a97",
  "source_dirty": true,
  "part": "xc7z020clg400-1",
  "top": "pynq_puf3_wrapper",
  "vivado": "2025.2",
  "control_id": "0x33465150",
  "hardware_validation": "NOT_RUN",
  "artifacts": {
    "pynq_puf3.bit": {
      "bytes": 4045693,
      "sha256": "769ed24c4ef4b155eccf5b92569a94d00315f7f6913b68c7eaffe65a96fa4654"
    },
    "pynq_puf3.hwh": {
      "bytes": 158472,
      "sha256": "4a0e7b95b6cb2d25c0a387e3a193aae56e937f0e2201346c67320dbc973a9675"
    },
    "pynq_svc_puf.bin": {
      "bytes": 47432,
      "sha256": "857df29c4310cd483e92730d81f333a68e6970658d5cb81344faf29ec50dd9b7"
    },
    "pynq_svc_puf.coe": {
      "bytes": 327742,
      "sha256": "bf3b0591afb251695d81846e6f844b2c5087b6ed38047242dcdde3fea05c6a9e"
    },
    "pynq_svc_puf.elf": {
      "bytes": 53868,
      "sha256": "d7ceb168daba6de16e67f8ef753b373ee0b87e2797ab2746297a45b5d3f3ff28"
    }
  }
}
\end{codeblock}


\Needspace{85pt}
\section{当前实现原件}
以下均来自 \pufreports，摘录行保留原文，省略的行不补写。
\subsection{时序原文}
\begin{codeblock}
| Tool Version : Vivado v.2025.2 (win64) Build 6299465 Fri Nov 14 19:35:11 GMT 2025
| Date         : Fri Oct 30 19:32:48 2026
| Design       : pynq_puf3_wrapper
| Device       : 7z020-clg400
| Speed File   : -1  PRODUCTION 1.12 2019-11-22
| Design State : Physopt postRoute
TIMING-17  Critical Warning  Non-clocked sequential cell                  1000        
| Design Timing Summary
    WNS(ns)      TNS(ns)  TNS Failing Endpoints  TNS Total Endpoints      WHS(ns)      THS(ns)  THS Failing Endpoints  THS Total Endpoints     WPWS(ns)     TPWS(ns)  TPWS Failing Endpoints  TPWS Total Endpoints  
      0.104        0.000                      0                73907        0.041        0.000                      0                73907        0.000        0.000                       0                 28016  
clk_fpga_0  {0.000 5.000}      10.000          100.000         
  ro_clk    {0.000 20.000}     40.000          25.000          
clk_fpga_0          0.104        0.000                      0                54424        0.041        0.000                      0                54424        0.000        0.000                       0                 22191  
  ro_clk           24.822        0.000                      0                 6127        0.115        0.000                      0                 6127       19.500        0.000                       0                  5825  
ro_clk        clk_fpga_0          3.742        0.000                      0                  258        0.216        0.000                      0                  258  
clk_fpga_0    ro_clk              6.994        0.000                      0                    1        0.428        0.000                      0                    1
\end{codeblock}

\subsection{整机资源原文}
\begin{codeblock}
| Tool Version : Vivado v.2025.2 (win64) Build 6299465 Fri Nov 14 19:35:11 GMT 2025
| Date         : Fri Oct 30 19:33:12 2026
| Design       : pynq_puf3_wrapper
| Device       : xc7z020clg400-1
| Design State : Physopt postRoute
| Slice LUTs                 | 18674 |     0 |          0 |     53200 | 35.10 |
| Slice Registers            | 32002 |     0 |          0 |    106400 | 30.08 |
| Slice Registers                            | 32002 |     0 |          0 |    106400 | 30.08 |
| Block RAM Tile    | 34.5 |     0 |          0 |       140 | 24.64 |
| DSPs           |    4 |     0 |          0 |       220 |  1.82 |
\end{codeblock}

\subsection{DRC 汇总原文}
\begin{codeblock}
| Tool Version : Vivado v.2025.2 (win64) Build 6299465 Fri Nov 14 19:35:11 GMT 2025
| Date         : Fri Oct 30 19:33:02 2026
| Host         : DESKTOP-5H4M0DR running 64-bit major release  (build 9200)
| Command      : report_drc -file C:/wenjian/game/FPGA2026/ours/amd_fpga_2026/build/local/pynq_puf3/vivado/reports/drc.rpt
| Design       : pynq_puf3_wrapper
| Device       : xc7z020clg400-1
| Speed File   : -1
| Design State : Fully Routed
             Checks found: 824
| Rule        | Severity | Description                                                 | Checks |
| CHECK-3     | Warning  | Report rule limit reached                                   | 2      |
| DPOP-1      | Warning  | PREG Output pipelining                                      | 2      |
| DPOP-2      | Warning  | MREG Output pipelining                                      | 4      |
| LUTLP-2     | Warning  | Combinatorial Loop Allowed                                  | 256    |
| PDRC-153    | Warning  | Gated clock check                                           | 256    |
| PLHOLDVIO-2 | Warning  | Non-Optimal connections which could lead to hold violations | 256    |
| RBOR-1      | Warning  | RAMB output registers                                       | 3      |
| REQP-1839   | Warning  | RAMB36 async control check                                  | 20     |
| REQP-1840   | Warning  | RAMB18 async control check                                  | 20     |
| RTSTAT-10   | Warning  | No routable loads                                           | 1      |
| UTLZ-3      | Warning  | Resource utilization                                        | 4      |
\end{codeblock}

\subsection{RO 计数器角分析与覆盖}
正 WNS 不覆盖未声明的振荡环计数器。下列器件模型分析不是板上 RO 频率测量。
\begin{codeblock}
There are 5120 register/latch pins with no clock driven by root clock pin: pynq_puf3_i/pynq_soc_0/inst/u_soc_axi_top/u_soc_top/u_puf_key/u_puf/ro_en_reg/Q (HIGH)
 There are 15360 pins that are not constrained for maximum delay. (HIGH)
 There are 256 combinational loops in the design. (HIGH)
FAST_MIN ring period ns: min 1.194 max 7.849 -> fastest RO 837 MHz
SLOW_MAX ring period ns: min 3.473 max 18.613 -> fastest RO 288 MHz
20-bit counter wraps above 400 MHz (window 65536 x 40 ns)
== Slow corner, RO period 3.473 ns ==
  ro_osc1 intra setup -0.478 (pynq_puf3_i/pynq_soc_0/inst/u_soc_axi_top/u_soc_top/u_puf_key/u_puf/g_ro[169].cnt_reg[169][3]/C -> pynq_puf3_i/pynq_soc_0/inst/u_soc_axi_top/u_soc_top/u_puf_key/u_puf/g_ro[169].cnt_reg[169][17]/D, skew -1.545)
  ro_osc1 intra hold  -0.552 (pynq_puf3_i/pynq_soc_0/inst/u_soc_axi_top/u_soc_top/u_puf_key/u_puf/g_ro[170].cnt_reg[170][15]/C -> pynq_puf3_i/pynq_soc_0/inst/u_soc_axi_top/u_soc_top/u_puf_key/u_puf/g_ro[170].cnt_reg[170][16]/D, skew 1.447)
== Fast corner, RO period 1.194 ns ==
  ro_osc1 intra setup -0.598 (pynq_puf3_i/pynq_soc_0/inst/u_soc_axi_top/u_soc_top/u_puf_key/u_puf/g_ro[182].cnt_reg[182][3]/C -> pynq_puf3_i/pynq_soc_0/inst/u_soc_axi_top/u_soc_top/u_puf_key/u_puf/g_ro[182].cnt_reg[182][11]/D, skew -0.705)
  ro_osc1 intra hold  -0.427 (pynq_puf3_i/pynq_soc_0/inst/u_soc_axi_top/u_soc_top/u_puf_key/u_puf/g_ro[170].cnt_reg[170][15]/C -> pynq_puf3_i/pynq_soc_0/inst/u_soc_axi_top/u_soc_top/u_puf_key/u_puf/g_ro[170].cnt_reg[170][16]/D, skew 0.804)
\end{codeblock}


\Needspace{85pt}
\section{容量内 KAT 核对}
来源：\pufresult 的 sections.kat。核对 3 组用例集合均为 PT/AD 各 0--16 B 的 289 组；KAT 编号 \#529 是其中一组，不是用例总数。
\begin{codeblock}
{
  "sw": {
    "kat_pass": 289,
    "tag_reject": 289,
    "case_count": 289
  },
  "hw": {
    "kat_pass": 289,
    "tag_reject": 289,
    "case_count": 289
  },
  "ise": {
    "kat_pass": 289,
    "tag_reject": 289,
    "case_count": 289
  }
}
\end{codeblock}


\Needspace{85pt}
\section{固定输入 warm 复算（KAT \#529，PT16/AD0）}
来源：\pufresult 的 sections.benchmark。5 次预热、每档 30 次；cycles 从 raw\_cycles 重算，墙钟从 wall\_ns\_summary 换算，二者证据粒度不同。
\begin{codeblock}
sw/encrypt: n=30 median_cycles=21862; cycles/B=1366.3750; median_wall_ms(summary)=5.1296550
sw/decrypt: n=30 median_cycles=22562; cycles/B=1410.1250; median_wall_ms(summary)=4.8743495
hw/encrypt: n=30 median_cycles=1369; cycles/B=85.5625; median_wall_ms(summary)=4.4026330
hw/decrypt: n=30 median_cycles=1213; cycles/B=75.8125; median_wall_ms(summary)=4.1595980
ise/encrypt: n=30 median_cycles=3895; cycles/B=243.4375; median_wall_ms(summary)=4.3854980
ise/decrypt: n=30 median_cycles=4288; cycles/B=268.0000; median_wall_ms(summary)=4.1371100
hw/encrypt speedup=15.969
hw/decrypt speedup=18.600
ise/encrypt speedup=5.613
ise/decrypt speedup=5.262
\end{codeblock}


\Needspace{85pt}
\section{旧 release 错误门禁记录}
来源：\texttt{data/pynq/error\_gate\_20261031\_112812/\allowbreak{}pynq\_board\_error\_gate.json}。这是旧 release 的证据，PUF3 完整错误矩阵仍为 NOT\_RUN。
\begin{codeblock}
firmware_id=pynq-svc-v1 (NOT PUF3)
initial_health: PASS
over_capacity_data_17B: PASS
over_capacity_ad_17B: PASS
half_frame_requires_reset: PASS
deadline_timeout_requires_recovery: PASS
duplicate_request_id_rejected: PASS
stale_epoch_rejected: PASS
rx_fifo_overflow_detected: PASS
repeated_requests: PASS
\end{codeblock}


\Needspace{85pt}
\section{归档仿真回归}
\subsection{PUF3 联合仿真与固件}
来源：\pufreports 的 cosim\_summary.json、firmware\_build.json。normal 为生产正例；mutant\_autolock 是应失败的原自动锁定 RTL 反例。仿真中 RO/DNA/BUFG 与 AXI BRAM 使用行为模型，不测 PUF 物理特性。
\begin{codeblock}
{
  "cosim": {
    "result": "PASS",
    "steps": 54,
    "normal_pass": true,
    "mutant_autolock_detected": true,
    "simulation_only": "behavioral RO/DNA/BUFG and AXI BRAM; KEY_WINDOW=256 via defparam"
  },
  "custom0": {
    "p12_0x200b": 1,
    "p8_0x300b": 4,
    "all_opcode_0x0b": 22
  },
  "size": "text\t   data\t    bss\t    dec\t    hex\tfilename\r\n  47424\t      4\t    344\t  47772\t   ba9c\tC:\\wenjian\\game\\FPGA2026\\ours\\amd_fpga_2026\\build\\local\\pynq_svc_puf\\pynq_svc_puf.elf"
}
\end{codeblock}

\subsection{ISE 截图配套单元仿真}
来源：\texttt{data/ascon/results/\allowbreak{}ise\_screenshot\_20261031/}。
截图显示首个 P8 片段；完整自检结果和源文件哈希另存，控制台的查询命令错误不计为 TB 断言错误。
\begin{codeblock}
Time resolution is 1 ps
SUMMARY ISE checks=1502 errors=0
$finish called at time : 37390 ns : File "C:/wenjian/game/FPGA2026/ours/amd_fpga_2026/sim/ascon/tb_key_ise.v" Line 63

{
  "status": "PASS",
  "evidence_kind": "behavioral_rtl_simulation",
  "tool": "Vivado/XSim 2025.2",
  "simulation_top": "tb_key_ise",
  "checks": 1502,
  "errors": 0,
  "finish_ns": 37390.0,
  "scope": "Existing ISE unit TB; not whole-SoC, timing, CDC, or board validation"
}
\end{codeblock}

\subsection{旧 P0 smoke 与 regression}
来源：\texttt{build/reports/pynq\_final\_20261028/}，保留 NOT\_RUN 项，不称为 PUF3 重跑。
\begin{codeblock}
{
  "smoke": {
    "source_head": "0f5d49c38597947dace3f1c99220deca8ef48e14",
    "counts": {
      "PASS": 6,
      "FAIL": 0,
      "BLOCKED": 0,
      "NOT_RUN": 9
    },
    "elapsed_s": 18.739,
    "status": "PASS"
  },
  "regression_final": {
    "source_head": "0f5d49c38597947dace3f1c99220deca8ef48e14",
    "counts": {
      "PASS": 15,
      "FAIL": 0,
      "BLOCKED": 0,
      "NOT_RUN": 1
    },
    "elapsed_s": 400.168,
    "status": "PASS"
  }
}
\end{codeblock}


\Needspace{85pt}
\section{host 标准验证与失败契约}
来源：\texttt{data/ascon/results/}。原始失败复现与生产修复后结果分别保留。
\begin{codeblock}
{
  "host_validation": {
    "passed": 28661,
    "failed": 0,
    "exit_code": 0
  },
  "reference_gate": {
    "passed": 12936,
    "failed": 0,
    "exit_code": 0
  },
  "failure_contract": {
    "passed": 23965,
    "failed": 0,
    "exit_code": 0
  },
  "failure_contract_original_reproduction": {
    "passed": 1,
    "failed": 1,
    "exit_code": 1
  }
}
\end{codeblock}


\Needspace{85pt}
\section{历史 PL-only 日志选行}
来源：\texttt{data/ascon/results/board/} 的 SW/ISE 和 MMIO 原始文本。当前性能主表不用这些历史周期。
\begin{codeblock}
=== ASCON-AEAD128 software on RV32IM ===
KAT empty : PASS
KAT 1B : PASS
KAT 16B : PASS
KAT 16B+AD : PASS
KAT 17B : PASS
KAT 32B : PASS
tamper reject : PASS
KAT: 6/6 PASS
enc 16B : 20806 cycles, 4495 inst, CPI 462 (x100)
enc 32B : 25787 cycles, 5606 inst, CPI 459 (x100)
dec 16B : 21360 cycles, 4722 inst, CPI 452 (x100)
dec 32B : 26444 cycles, 5857 inst, CPI 451 (x100)
===DONE===
=== ASCON-AEAD128 ISE (permutation in hardware) ===
KAT empty  : PASS
KAT 1B     : PASS
KAT 16B    : PASS
KAT 16B+AD : PASS
dec round  : PASS
tamper     : PASS
ISE cycles   : 2570
===DONE===
=== ASCON-AEAD128 hardware (MMIO) on RV32IM ===
KAT empty   : PASS
KAT 1B      : PASS
KAT 16B     : PASS
KAT 16B+AD  : PASS
dec roundtrip: PASS
tamper reject: PASS
len reject  : PASS
crypto latency : 103 cycles
full MMIO      : 688 cycles
===DONE===
\end{codeblock}


\Needspace{85pt}
\section{历史 PL-only 时序原文}
来源：\texttt{build/local/\allowbreak{}proj\_backup\_first\_time\_20261027/\allowbreak{}first\_time.runs/\allowbreak{}impl\_1/} 的 routed timing summary。本地原件本次可读，但不保证包含在远端 clone 中。
\begin{codeblock}
| Tool Version : Vivado v.2025.2 (win64) Build 6299465 Fri Nov 14 19:35:11 GMT 2025
| Date         : Tue Oct 27 18:42:44 2026
| Design       : top_pl_bram
| Device       : 7z020-clg400
| Speed File   : -1  PRODUCTION 1.12 2019-11-22
| Design State : Routed
| Design Timing Summary
    WNS(ns)      TNS(ns)  TNS Failing Endpoints  TNS Total Endpoints      WHS(ns)      THS(ns)  THS Failing Endpoints  THS Total Endpoints     WPWS(ns)     TPWS(ns)  TPWS Failing Endpoints  TPWS Total Endpoints  
      0.087        0.000                      0                55572        0.021        0.000                      0                55572        3.750        0.000                       0                 17227
\end{codeblock}


\Needspace{85pt}
\section{历史策略 sweep 记录}
逐项记录见 \texttt{strategy\_sweep\_summary.json}，位于 \texttt{data/ascon/results/\allowbreak{}strategy\_sweep\_20261017/}。该历史默认/最优 WNS 为 $-0.942/-0.118$ ns；不与不同源集的 PUF3 run 混作单补丁效果。校验清单包含该文件哈希。

\Needspace{85pt}
\section{实物照片与工程截图}
图 3-1 为已有板卡正面照片；图 2-10、2-11、3-2、3-5 分别为用户提供的 Block Design、Implemented Design、Utilization 和 ISE 波形截图，每张只出现一次。四张 PNG 逐字节复制并保存 SHA-256，工程名、数值、波形及控制台记录保持原样。
\begin{codeblock}
vivado_block_design.png
SHA-256 0bd4c00d440e6383c67189dfdfd83e5fc6279eb30ec67ee4a6b540aa3e9201c6
vivado_implemented_design.png
SHA-256 7931afa35924144401781e4f1529f6bca8d9d5b585a5a7dd0c3b52d96c023e91
vivado_utilization.png
SHA-256 1047372c228582332d7d1c4abcd6cb79dcfa2d56fffb9dec91173f2c4952b2cf
vivado_ise_waveform.png
SHA-256 4fef2cc0498695c3c9d8a11f3ad4bfd0924e2433d8f3914090eee160f06ad4ac
\end{codeblock}


\Needspace{85pt}
\section{PUF3 功能与历史原型边界}
当前来源：\pufresult。下列结果证明会话内功能，不证明跨上电恢复、BER、熵、温压或唯一性；REF/RAW/MASK/DNA 可读，也不证明保密隔离。
\begin{codeblock}
{
  "clock": {
    "fclk0_mhz_register_derived": 99.999999,
    "fclk0_mhz_measured_mcycle": 99.9723,
    "pynq_clocks_fclk0_mhz": 150.0,
    "measurement": {
      "mcycle": [
        1417316,
        302040308
      ],
      "host_s": [
        353.655703317,
        356.662766276
      ]
    }
  },
  "functional": {
    "enroll_quality": 127,
    "reconstructs_stable": 5,
    "caller_key_rejected_under_puf_key": true,
    "same_ciphertext_after_reconstruct": true,
    "kat_restored_after_hwkey_off": true,
    "puf_key_hw_encrypt_cycles": 1676,
    "limits": "in-session only (helper volatile); REF/MASK/RAW/DNA readable; entropy/BER/temperature/cross-power-cycle not measured"
  },
  "timing_signoff_ro_domain": false,
  "cross_power_cycle": "not_measured"
}
\end{codeblock}

历史来源：\texttt{data/ascon/results/board/uart\_20261027\_keys3.txt}，保留 v2 抢跑失败和 v3 通过片段；历史 quality=128 不覆盖当前的 127。
\begin{codeblock}
# 来源：XCOM 串口捕获（COM7 / 115200 8N1），原始字节见 xcom_raw_20261027_keys3.XDat
# 本文件含两段连续捕获：
#   第 1 段 = PUF 握手 v2 固件（固件未等 idle 抢跑，读数误判，结论作废，仅作问题记录）
#   第 2 段 = v3 固件（先等 idle 再发 START），板级全通过
[第 1 段：v2 固件 — 抢跑误判]
STATUS boot   : 00000001 busy=1 enrolled=0 done=0 quality=0
STATUS pre-enr: 00000001 busy=1 enrolled=0 done=0 quality=0
STATUS enrol: 00000023 quality=0
FAIL enroll not marked
FAIL quality=0
STATUS pre-rec: 00007f05 busy=1 enrolled=1 done=0 quality=127
STATUS recon: 00007f26
FAIL
[第 2 段：v3 固件 — 板级通过]
STATUS boot   : 00000001 busy=1 enrolled=0 done=0 quality=0
STATUS idle   : 0000003a busy=0 enrolled=0 done=1 quality=0
STATUS enrol: 00008026 quality=128
STATUS recon: 00008026
PASS key subsystem: enroll + reconstruct stable
\end{codeblock}

